\documentclass[11pt]{article}
\usepackage{amsmath,amssymb,amsthm}
\usepackage[margin=1.05in]{geometry}
\usepackage[hidelinks]{hyperref}
\newtheorem{theorem}{Theorem}
\newtheorem{lemma}[theorem]{Lemma}
\newtheorem{proposition}[theorem]{Proposition}
\newtheorem{corollary}[theorem]{Corollary}
\title{Erd\H{o}s Problem \#982: Improved Bounds, Computer-Assisted Proofs, and Structural Results}
\author{}
\date{26 November 2026 \\ Repository manuscript for review}
\begin{document}
\maketitle
\begin{abstract}
Erd\H{o}s asked whether every planar strictly convex \(n\)-gon has a vertex at at least \(\lfloor n/2\rfloor\) different distances from the other vertices. We give a lower bound of \((13/36+1/1667)n-O(1)\) by controlling overlapping windows in an iteration of enclosing-disk supports. A separating tree yields \(\lceil n/2\rceil\) for samples with no collinear triples and rational coordinates up to similarity, and for real number fields with a specified dyadic ramification property. We give the finite reduction used with computer refutations for \(3\le n\le19\) and \(n=21\), and prove results for cocircular samples, conics, constant-width boundaries, and metric perturbations. The unrestricted conjecture and the exact twenty-vertex value remain open in the audited material.
\end{abstract}

The largest distance count from one pin and the number of distinct distances over all pairs are different invariants. We study the first, as in Erd\H{o}s's problem and its subsequent lower-bound literature \cite{erdos,moser,ef,nppz,kominers}. Strict convex position means that every point is a vertex of the convex hull. In particular there are no three collinear vertices. All planar distances are Euclidean unless a metric is specified. Positive distances exclude the pin itself.

The observation of no collinear triples in the higher algebraic-coordinate bound is needed in the pair count. The finite reductions use strict convexity to impose additional ordering and angle constraints. The certificate catalogue accompanying this manuscript records coverage, exact bytes and which checker was used for each branch. Contributions and literature overlap are recorded by stable identifiers in the repository. Proof material under review is indexed there as well.

\section{Minimum enclosing disks and finite reductions}

Let \(f(n)\) be the minimum of \(D(P)\) over the vertex sets of strictly convex \(n\)-gons. For any subset \(S\subset P\), \(D(P)\ge D(S)\). The regular \(n\)-gon has \(D(P)=\lfloor n/2\rfloor\).

\begin{lemma}[Support and monotonicity]
The minimum enclosing disk of \(P\) is supported either by a contact diameter or by three contact vertices forming an acute triangle. On each polygonal chain between successive support vertices, distances from any vertex strictly increase along either direction of traversal away from it.
\end{lemma}

\begin{proof}
If its center did not lie in the convex hull of the contact vertices, strict separation would give a common direction of strictly decreasing distance to every contact vertex. The finitely many noncontact vertices have positive slack. A sufficiently small move would contradict minimality. By planar Carath\'eodory, the center belongs to the convex hull of at most three contact points. A two-point support is a diameter. With three points and no such diameter, the center is interior to their triangle, which is therefore acute.

For a support side \(AB\) let \(\gamma\) be the angle at the third contact point. Along the polygonal chain facing the minor circular arc, each interior vertex \(v\) satisfies \(\angle AvB\ge\pi-\gamma>\pi/2\) by the inscribed-angle comparison in the closed cap. The diameter case gives \(\angle AvB\ge\pi/2\) on either chain. If \(x,v,y\) occur in this order on the same chain, the rays to \(A,B\) lie in \(\angle xvy\) by convexity. Thus \(\angle xvy\ge\pi/2\), and the side \(xy\) is strictly longer than \(xv\) and \(vy\). Applying this to any three such vertices proves the monotonicity.
\end{proof}

If \(P\) has three vertices, each pin sees at least one positive distance. If it has four, the support lemma gives a chain with at least two edges. Deleting one vertex extends the four-vertex result to five vertices.

For an even size \(2h\), a counterexample has no contact diameter and all three support gaps are at most \(h-1\). The gaps count edges and sum to \(2h\). Dihedral changes of the boundary labels realize any permutation of the three gap positions. At \(n=20\) the resulting unordered triples are
\[
 (9,9,2), (9,8,3), (9,7,4), (9,6,5), (8,8,4),
 (8,7,5), (8,6,6), (7,7,6).
\]
In particular the diameter condition concerns antipodal contacts of the minimum disk. For a diameter pair in an arbitrary point set, containment in the disk having that pair as a diameter requires a separate argument.

\paragraph{From geometry to clauses.}
A putative counterexample induces local distance ranks, numbered from zero. In the threshold encoding \(T_{p,q,k}\) means the rank is at least \(k\). The comparison literal \(G_{p,a,b}\) means \(pa>pb\). Its truth implies \(T_{p,b,k}\Rightarrow T_{p,a,k+1}\); conversely \(T_{p,a,k}\) with \(T_{p,b,k}\) false forces the comparison. The global-order encoding instead compares the real lengths of all unoriented edges. In either encoding, real lengths cannot increase weakly around a cycle with one strict step. Diagonals of a convex quadrilateral have sum greater than the sum of either opposite pair, by splitting diagonals at their interior intersection and applying two strict triangle inequalities. Each support-chain clause follows from the support lemma.

At a vertex the total angle between the extreme rays is strictly below \(\pi\). Three disjoint sectors whose opposite edges are longest in the corresponding triangles would have angles summing to at least \(\pi\). At an interior cap vertex the two sectors outside the endpoint rays have total angle strictly below the opposite support angle, hence below \(\pi/2\). There cannot be two disjoint sectors of size at least \(\pi/3\) there. For generators with angle bins, use \(\lfloor 3\theta/\pi\rfloor\) and the predicate \(\theta\ge\pi/2\). A sum of two consecutive angles has bin equal to the sum of bins or one greater. The bins of a nondegenerate triangle sum to \(1\), \(2\), or \(3\). Angles compared in a single triangle are ordered by their opposite sides. For a cyclically ordered quadrilateral \(i,j,v,k\), the premise \(iv\le ij,ik\) implies that the angle at \(v\) is at least the sum of the angles at \(j,k\), by applying the triangle side-angle comparison on each side of \(iv\). Consequently \(2\angle jvk\ge2\pi-\angle jik\). A shortest support side has opposite angle at most \(\pi/3\), which strengthens the cap angle and the exterior budget by the support lemma.

The refinements using the entire enclosing disk and the enumeration of weak orders and missing ranks are recorded in the finite encoding audit accompanying this manuscript. They are applied only in the generator families identified there. An empty clause derived by reverse unit propagation from a subset of a rebuilt formula excludes every real counterexample satisfying the premises of the corresponding branch.

\paragraph{Twenty-one vertices.}
The support lemma settles a contact diameter and any chain of at least ten edges. If the longest gap is nine, delete a noncontact vertex from one of the other two chains with an interior vertex. The acute contact triangle still forces the same minimum disk. The remaining gaps have sum twenty and maximum nine and hence are one of the four nine-leading triples above. If all gaps are at most eight on twenty-one vertices, their only possibilities are \((8,8,5),(8,7,6),(7,7,7)\). Thus exclusion of these three profiles and the four specified twenty-vertex profiles suffices for \(f(21)=10\). For odd \(n=2h+1\) more generally, any exact bound \(f(2h)=h\) implies \(f(2h+1)=h\). This implication uses deletion from the odd size to the even size.

\section{Cocircular equality and symmetric samples}

\begin{theorem}[982-C01, 982-C02, 982-C03]
For \(m\) distinct points on a circle every pin sees at least \(\lfloor m/2\rfloor\) distances. Simultaneous equality for odd \(m\ge3\) holds precisely for the regular \(m\)-gon. For even \(m\ge6\), it holds precisely for the regular \(m\)-gon, the regular \((m+1)\)-gon with a vertex removed, or the regular \((m-1)\)-gon completed by the antipode of one of its vertices.
\end{theorem}

\begin{proof}
Normalize the circle to the complex unit circle. At \(p\), the only possible pairing of distinct targets at equal distance is by \(R_p(x)=z\bar x\), with \(z=p^2\). If this reflection has \(f_p\) fixed sample points and \(u_p\) unmatched sample points, counting pairs and then excluding distance zero gives \(d_P(p)=(m+f_p+u_p)/2-1\), with \(f_p\ge1\). This proves the lower bound. In the odd equality case, all radial reflections preserve the set. Each such reflection equates the adjacent circular gaps at that vertex, so the set is regular.

In the even equality case \(u_p+f_p=2\). Put \(S=\sum_{x\in P}x\) and \(T=\sum_{x\in P}x^2\). If \(u_p=1\) and \(y\) is the unmatched point, then
\[
 S-z\bar S=y-z/y,\qquad
 T-z^2\bar T=y^2-z^2/y^2.
\]
If \(u_p=0\), both left sides vanish. In both cases \(z\) is a root of
\[
 (T-z^2\bar T)^2-(S-z\bar S)^4-4z(S-z\bar S)^2.
\]
If \(S=0\) no point can be unmatched and all radial reflections preserve the set. Otherwise every radial symmetry axis passes through the nonzero centroid. There is at most one such axis, and every antipodal pair lies on one, so there is at most one antipodal pair. The polynomial has at least \(m-1\ge5\) distinct roots and vanishes identically. Rotate so \(S=s>0\). The constant coefficient gives \(T=\pm s^2\). Substitution gives \(4z(1-z)^2s^2(s^2-1)=0\) identically, hence \(s=1\).

For \(z\ne1\), the two moment identities give \(y=1\) if \(T=1\) and \(y=-z\) if \(T=-1\). If \(T=1\), remove \(1\); the remaining set is invariant under the radial reflections of its points, including any real-axis pin. It is a regular odd polygon, and invariance under the real-axis reflection forces \(-1\) to be a vertex. If \(T=-1\), add the missing \(-1\); this restores invariance under every reflection based at an old vertex, including any real-axis pin. The adjacent-gap equalities at all but the added vertex already imply equal gaps. Conversely, deletion from a regular odd polygon retains at least one member of each paired distance. Adding an antipode to a regular odd polygon adds at most one distance at each old pin, and the added pin sees one singleton and the remaining targets in pairs. The regular even polygon also attains the bound.
\end{proof}

\begin{theorem}[982-C04, 982-C05]
A strictly convex even \(n\)-gon with \(n\ge8\) satisfies \(D(P)\ge n/2\) if at least \(n-2\) vertices are cocircular. For odd \(n\ge9\), at least \(n-3\) cocircular vertices suffice.
\end{theorem}

\begin{proof}
If the circle center is a vertex, a supporting line there puts the cocircular vertices in an open semicircle. From the first of them all chord lengths increase, giving at least \(n-3\) distances in the even case. With all vertices on the circle the preceding theorem applies. With exactly one external vertex other than the center, a circle about that vertex intersects the original circle at most twice at each radius, giving \(n/2\) distances.

For exactly two off-circle vertices \(q,r\), put \(m=n-2\). If every pin fails the desired bound, the on-circle set \(A\) has simultaneous equality and each of \(q,r\) sees the \(m\) on-circle targets perfectly paired. Thus their radial reflections preserve \(A\) and have no fixed vertex on it.
If \(A\) is a regular \(m\)-gon, write \(\alpha=\pi/m\) and let \(\rho=|q|>0\) on the unit circle normalization. Its radial axis bisects opposite edges. All \(m/2\) distances from \(q\) to \(A\) must equal the old chord lengths as an ordered list, since every old pin already uses its full distance budget. This requires
\[
 \rho^2-1+2(\cos\alpha-\rho)\cos x-2\sin\alpha\sin x=0
\]
at \(x=(2j-1)\alpha\), \(1\le j\le m/2\). A nontrivial line has at most two intersections with the unit circle, whereas \(m/2\ge3\).

If \(A\) lacks vertex \(1\) of a regular \((m+1)\)-gon, its vector sum is \(-1\). A symmetry axis fixes the centroid, so \(q,r\) lie on the real axis. Each must see the original \(m/2\) chord lengths twice, exactly as pin \(1\) would. Summing their squares for either real pin \(t\) gives \(m(1+t^2)+2t=2m+2\). The only off-circle root is \(-1-2/m\), contradicting distinctness. In the remaining classification the vector sum is the added point, so the only possible radial symmetry axis contains an antipodal pair and has fixed vertices. This finishes the even case. Delete an off-circle vertex if possible for the odd case and apply the even result.
\end{proof}

\begin{theorem}[982-S01]
Let \(P\) be a finite sample of the boundary of a centrally symmetric strictly convex planar body with center \(O\). Each pin \(A\in P\) farthest from \(O\) sees each positive distance at most twice.
\end{theorem}

\begin{proof}
Adjoin the antipode \(B=2O-A\) if necessary. All sample points lie in the disk centered at \(O\) of radius \(OA\). The enlarged sample is still in strictly convex position because it lies on the same strictly convex boundary. The disk with diameter \(AB\) is its minimum enclosing disk, so the two chains have strictly increasing distance from \(A\). There are at most two occurrences across the chains, and deletion of the added point preserves this bound.
\end{proof}

\begin{theorem}[982-S02]
For a strictly convex planar sample from a nondegenerate real conic, some pin sees every positive distance at most twice.
\end{theorem}

\begin{proof}
Ellipses are covered by 982-S01. For a sample on the graph of a convex function, take a pin of minimum sample height. To its right the first sampled secant has nonnegative slope, and all subsequent adjacent secants have at least this slope by convexity. Thus the height is nondecreasing and distance from the pin is strictly increasing. Reflecting the argument handles the left. A parabola can be put in this form by an isometry, as can either branch of a hyperbola, after reflecting a concave branch if necessary.

If a strictly convex sample meets both branches of a hyperbola \(y=\pm b\sqrt{1+x^2/a^2}\), each branch contains at most two samples. In fact, if three upper samples were present, a strict supporting line of the middle one \(q=(x_0,y_0)\) would have to lie below the two others, with slope strictly between their two secant slopes from \(q\). By the derivative bound this slope \(k\) satisfies \(|k|<b/a\), and \(y_0-kx_0>0\). The line then lies strictly above the entire lower branch, contradicting strict support. The reflected argument bounds the other branch. At most four samples remain, and any non-equilateral triangle among them gives a pin with at least two distinct distances. Four distinct planar points cannot be mutually equidistant, so a qualifying pin always exists.
\end{proof}

\section{An improved unrestricted bound}

\begin{theorem}[982-G05]
There is an absolute constant \(C\) such that every strictly convex \(n\)-gon satisfies
\[
 D(P)\ge\left(\frac{13}{36}+\frac1{1667}\right)n-C.
\]
\end{theorem}

\begin{proof}
Call a pair good if its perpendicular bisector contains at most one vertex, and write \(G\) for the number of good pairs. The number of isosceles triples counted by their apex is at most \(2\binom n2-G\), since a bisector is a line. Counting at each apex and using \(\binom r2\ge2r-3\) for a nonempty class of size \(r\) gives
\[
 D(P)\ge(n-1)/3+G/(3n).
\]
We prove \(G\ge(1/12+\eta)n^2-O(n)\), where \(\eta=3/1667\).

\paragraph{Cap counts.}
In a cap, every ordered triple \(u,x,v\) has middle angle at least \(\pi/2\). A pair of cap vertices is good if the intervening chain in the original polygon contains no bisector witness. Indeed, a diagonal has at most one boundary vertex on its bisector on each side. If \(p,q\) are left of a witness \(z\), and \(v\) is to its right, the order \(p,q,z,v\) shows \(\angle pqz\ge\angle pqv\ge\pi/2\) by convexity. Thus \(zp>zq\), even when the witness \(z\) has been removed and \(p,q,v\) are the remaining cap vertices. The symmetric statement holds on the right. It follows that, for each witness, the pairs witnessed within a cap are a matching between its two sides. For an intact cap of \(l\) vertices, the total of \(\min(j-1,l-j)\) over positions \(j\) therefore gives at least \(l^2/4-O(l)\) unwitnessed pairs. The same count applies to a residual cap before accounting for deleted witnesses.

For an endpoint pair \(uv\) of a cap, any witness of \(uc\), with \(c\ne v\), is closer to \(u\) than to \(v\), by the same quadrilateral argument with \(z\) between \(u,c\). Conversely a witness of \(cv\), \(c\ne u\), is closer to \(v\). A witness of \(uv\) cannot witness either other type. Hence each vertex in the original chain spoils at most one endpoint pair of this cap. Among \(2l-3\) endpoint pairs at least \(l-1\) lack a witness in the remaining cap. We also use Nivasch--Pach--Pinchasi--Zerbib's cap lemma \cite[Lemma 12]{nppz}: between the first and last \(t\) vertices of a cap on \(2t\) vertices there are at least \(t^2/8-O(t)\) pairs with no intervening witness in that cap. Deleted points outside its original chain cannot spoil these pairs.

\paragraph{Initial partition and parameters.}
A contact diameter gives two caps and \(G\ge n^2/8-O(n)\). Otherwise let \(w_1,w_2,w_3\) be the support triangle. Write \(\lambda_j\) for the three fractions of boundary edges on its chains, \(E=\frac14\sum_j(\lambda_j-1/3)^2\), \(\rho=\min_j\lambda_j=1/3-x\). The initial caps give \((1/12+E)n^2-O(n)\) good pairs, and \(E\ge3x^2/8\). If \(E\ge\eta\) we are done. Fix \(\delta=\eta-3x^2/8>0\) in the remaining case. Boundary distances below refer to the original cycle, divided by \(n\).

Let \(d=353/40000\). In each of \(M=\lfloor dn\rfloor\) rounds remove a support set of the current minimum disk. Write \(s=i/n\) in round \(i\), and let \(r_i(w)\) denote the shortest boundary distance from \(w\) to a support vertex in that round. Set
\[
 a_z(s)=5s+\sqrt{16s^2+8z},\qquad
 \psi_s(r)=
 \begin{cases}
 0 & r\le9s,\\
 (r-s)(r-9s)/8 & r>9s.
 \end{cases}
\]
Here \(\psi_s\) is defined also for negative arguments, is increasing on \([9s,\infty)\), and is convex. Let
\[
 \begin{split}
 U(s)&=a_\delta(s)+s,\quad c(s)=\rho-U(s),\\
 v&=(r-c(s))_+,\quad q=\min(s,v/6),\quad t=(r-s-q)_+,\\
 \Phi_s(r)&=(t^2/8-st-v^2/12)_+ .
 \end{split}
\]
The following parameter inequalities hold for \(0\le s\le d\) and \(0<\delta=\eta-3x^2/8\):
\begin{align}
 c(s)&>a_{3\delta/5}(s), & 2U(s)-\rho&<6s, \label{eq:params}\\
 \Phi_s(r)&\ge\psi_s(r-s) &&(c(s)\le r\le U(s)), \label{eq:shift}\\
 c(s)+s+2a_{(\delta-\psi_s(c(s)))/2}(s)&\le3a_{\eta/3}(s)
 &&(c(s)<U(s),\ \psi_s(c(s))\le\delta), \label{eq:sum}\\
 d/2-9d^2-\eta&>3d\sqrt{(16d^2+8\eta)/3}. \label{eq:integral}
\end{align}
Their exact verification is given below. In particular \(c(s)>9s\) and \(\Phi_s=\psi_s\) up to \(c(s)\). At fixed \(s\), \(\Phi_s\) is nondecreasing, and strictly increasing beyond \(c(s)\). For \(v\le6s\) beyond the cutoff, the derivative before truncation is \(5(c(s)-5s)/24+v/144\); for \(v>6s\) it is \((r+2c(s)-18s)/12\). Both are positive. At fixed \(r\), \(\Phi_s(r)\) is nonincreasing as \(s\) increases: \(v,q\) increase and \(t\) decreases, and on a later positive value \(t>8s\), where \(t^2/8-st\) increases with \(t\).

\paragraph{First round with a large sum.}
Stop at the first \(i\) for which \(\sum_w\Phi_s(r_i(w))>\delta\), if there is such a round. In every preceding round all \(r_j(w)\le U(j/n)\), since \(\Phi_s(U(s))\ge\delta\) by \eqref{eq:shift} or the no-overlap formula and is strictly increasing beyond \(U(s)\). No support vertex can be nearest to two centers: their two distances would sum to at least \(\rho\), each would be at least \(c(s)\), and their combined cost would exceed \(\delta\). Thus the three nearest contacts are distinct and account for all deletions of each preceding round. Round one has \(r_1(w)=0\).

In the stopping round choose radii \(0\le r_w\le\min(r_i(w),U(s))\) with total cost \(\delta\). If some distance exceeds \(U(s)\) the one radius suffices, by continuity and the bound at \(U(s)\); otherwise continuously shrink the three radii from their actual values. Ignore windows with zero cost. Any two chosen radii have sum strictly below \(\rho\), as otherwise each is at least \(c(s)\), contradicting the total cost. In particular each active window lies across its own center, contains no other original center and is disjoint from the others. It contains no current support vertex because the interval is open.

In the window of radius \(r\) about \(w\), at most \(sn\) earlier deletions matched \(w\). An earlier deletion matched to another center can lie inside only if its distance to its own center is greater than \(\rho-r\ge c(s)\). By monotonicity in time its cost at the earlier time exceeds \(3\delta/5\). Hence there is at most one such deletion per earlier round, even if the two outer windows overlap this one. From \(U'\ge6\), this intrusion can only occur during the last \(qn+O(1)\) earlier rounds. At time \(u=j/n\), the intruder lies within \(n(r+U(u)-\rho)_++O(1)\) steps of the outer edge of the window on its side, by the triangle inequality for cycle distance. Each spoiled cross pair uses a different surviving endpoint between that intruder and that outer edge, by the matching observation. The total cost for intruders is at most
\[
 n^2\int_0^s(r+U(u)-\rho)_+\,du+O(n)\le v^2 n^2/12+O(n).
\]
The sum-to-integral error is uniform, since \(6\le U'\le10\). On each side of \(w\) choose the \(tn+O(1)\) closest surviving vertices, with the same number on both sides. They occur consecutively in one current cap. The cross lemma gives \(t^2n^2/8-O(n)\) candidate pairs; the earlier deletions matched to \(w\) spoil at most \(stn^2+O(n)\) of these by the matching observation. Subtracting the intruder bound gives \(\Phi_s(r)n^2-O(n)\) good pairs for each active window. Their ends belong to distinct initial caps, and pairs from different windows are distinct. Adding the initial cap count proves the desired bound since \(E+\delta\ge\eta\).

\paragraph{All rounds have a small sum.}
In this case every round has exactly three matched support vertices. If all \(r_i(w)\le c(s)\), convexity of \(\psi_s\) gives \(\sum_w r_i(w)\le3a_{\eta/3}(s)\). Otherwise exactly one exceeds \(c(s)\). Set \(z=\Phi_s(r_i(w))>3\delta/5\) for that window. By \eqref{eq:shift} and convexity the sum is at most \(s+a_z(s)+2a_{(\delta-z)/2}(s)\). As a function of \(z\) this bound is nonincreasing for \(z\ge\delta/3\), by differentiation or concavity of the square root. Substitute \(z=\psi_s(c(s))\) and apply \eqref{eq:sum}.

In round \(i\) the three cap endpoint counts give at least \(n-3(i-1)\) pairs lacking witnesses in the remaining caps. Each earlier deletion spoils at most one such pair, so the rounds supply \((d-3d^2)n^2-O(n)\) good pairs, disjoint because at least one endpoint is deleted in the current round. In the remaining parts of the original caps the intact count gives \((n-3M)^2/12-O(n)\) further pairs before accounting for deleted witnesses. For a deleted vertex \(p\) matched to \(w\), a spoiled residual pair has one endpoint between \(p\) and \(w\) in their initial cap. Distinct spoiled pairs at \(p\) use distinct such remaining endpoints. If \(b(p)\) denotes the unnormalized boundary distance from \(p\) to \(w\), there are initially at most \(b(p)\) possible endpoints. For each center, the \(M-1\) noncentral deleted vertices are on its two sides. If their counts are \(\ell,r\), then \(\binom\ell2+\binom r2\ge (M-1)^2/4-(M-1)/2\) potential endpoints have already been deleted, counted once in the budget of each farther vertex they precede on their side. This gives total losses at most
\[
 3n^2\int_0^d a_{\eta/3}(s)\,ds-\tfrac34 d^2 n^2+O(n).
\]
The net coefficient beyond \(1/12\) is \(d/2-3d^2/2-3\int_0^d a_{\eta/3}(s)\,ds\). Concavity of the square root and averaging \(s^2\) bound the integral by \(5d^2/2+d\sqrt{(16d^2+8\eta)/3}\). Equation \eqref{eq:integral} now gives the result. All counts use only bounded numbers of caps or windows at a time and the displayed linear errors have constants independent of \(\delta\) and \(P\).
\end{proof}

\paragraph{Exact parameter verification.}
In \eqref{eq:params} the differences to check are worst at \(s=d\). Squaring positive sides and applying
\((\sqrt u+\sqrt v)^2\le9u/5+9v/4\) for the first inequality gives quadratic polynomials of \(x\) with positive leading coefficient and negative discriminant. For \eqref{eq:shift}, subtract the untruncated quadratics and multiply by \(288\); the difference is \(72rs-396s^2-12(r-5s)v-23v^2\). When the target \(\psi_s(r-s)\) vanishes there is nothing to prove. Otherwise the displayed polynomial is concave on \([c(s),U(s)]\). At \(c(s)\) it is nonnegative. At \(r=U(s)\), writing \(w=\sqrt{16s^2+8\delta}\), its derivative multiplied by \(w\) is
\[
 -4(4800s^2+2638sw+300w^2-416\rho s-141\rho w).
\]
In the overlapping regime \(\rho<12s+2w\) the parenthesis is positive by \(w\ge4s\). The minimum is therefore at \(s=d\). On \(x\in[3/50,7/100]\) a quadratic root count rules out overlap; the allowed \(x\) lie below \(7/100\). Below \(3/50\) the endpoint polynomial reduces to \(P(x)+Q(x)w\), where \(Q>0\) and a Sturm count proves \(Q^2w^2-P^2>0\).

For \eqref{eq:sum}, subdivide \(0\le s\le d,\ 0\le x\le3/50\) into \(384\) by \(768\) closed rational rectangles. In each rectangle bound \(U\) from above using the largest \(s,\delta\); skip it if this bound is at most half the lower endpoint bound for \(\rho\). In every other rectangle let \(L\) be the lower bound for \(c(s)\), \(s_+\) and \(s_-\) the upper and lower time endpoints, and \(\delta_+\) the upper bound for \(\delta\). Check \(L>a_{\delta_+/3}(s_+)\). For fixed \(s,\delta\), the map \(r\mapsto r+2a_{(\delta-\psi_s(r))/2}(s)\) decreases on \(r\ge a_{\delta/3}(s)\) where \(\psi_s(r)\le\delta\): its derivative equals \(1-(r-5s)/\sqrt{16s^2+4(\delta-\psi_s(r))}\le0\). Thus replacing \(c(s)\) by \(L\), then bounding the remaining square root with \(s_+\), \(\delta_+\) and \(z_-=(L-s_+)(L-9s_+)/8\), gives a valid upper bound. After subtracting \(11s\) from both sides the rational checks compare
\[
 L+2\sqrt{16s_+^2+4\max(0,\delta_+-z_-)}
 \quad\text{with}\quad 4s_-+3\sqrt{16s_-^2+8\eta/3}.
\]
Square-root endpoints are enclosed by integer square roots on the \(10^{-8}\) lattice. The accompanying script verifies all these comparisons and \eqref{eq:integral} by exact arithmetic; decimal values printed by that script are summaries of rational inequalities.

\section{Binary distance separation and coordinate fields}

Write \(d_P(p)=|\{|p-q|:q\in P\setminus\{p\}\}|\) and \(D(P)=\max_{p\in P}d_P(p)\).
A binary tree on \(P\) is separating if, at every split, any pin in that node has distinct distances to targets in different children, with the pin itself omitted.

\begin{theorem}[982-D01]
Let \(P\) consist of \(n\ge2\) distinct planar points with no three collinear. If \(P\) admits a separating binary tree, then \(D(P)\ge\lceil n/2\rceil\).
\end{theorem}

\begin{proof}
In a node \(S\) of size \(m\), let the two children have sizes \(b\le c=m-b\).
Count the distinct distances separately for each pin in \(S\) and each child. A target pair has at most two witnesses in \(S\), since its bisector is a line. For the \(b\)-child use \(r-1\le\binom r2\) on each nonempty distance class of size \(r\). The total contribution is at least \(b(m-1)-b(b-1)=bc\). For the \(c\)-child use \(2r-3\le\binom r2\), giving at least \((2c(m-1)-c(c-1))/3=c(m+b-1)/3\). Thus
\[
 3\sum_{p\in S}d_S(p)\ge(m-b)(m+4b-1).
\]
On the path to \(S\), let the discarded siblings have sizes \(t_j\) and set \(V=\sum_j t_j^2\). For pins in \(S\), each sibling contributes at least \(mt_j-t_j(t_j-1)\) by the same pair count. Distances in distinct sibling groups and inside \(S\) are disjoint at each pin by separation at the first split between the targets. Hence
\[
 3mD(P)\ge(m-b)(m+4b-1)+3(m+1)(n-m)-3V.
\]
At the root \(m=n\) and \(V=0\). Maintain \(m>n/2\) and \(3V\le m(n-m)\). If \(b<(2m-n)/4\), descend into the larger child. The new size \(m-b\) exceeds \(n/2\) and the slack in the second invariant increases by \(b(2m-n-4b)>0\). The positive integer size decreases, and a leaf cannot meet \(m>n/2\), so eventually \(b\ge(2m-n)/4\). At this stopping node
\[
 3m\bigl(D(P)+1/3-n/2\bigr)
 \ge m(n/2-m)+3mb-4b^2+b+3(n-m).
\]
The first three terms on the right are concave in \(b\), with values \((m-n/2)(n-m)/2\) at \((2m-n)/4\) and \(m(n-m)/2\) at \(m/2\). The right side is nonnegative. Since \(D(P)\) is an integer, \(D(P)\ge n/2-1/3\) implies the stated bound in both parities.
\end{proof}

\begin{theorem}[982-D02]
Let \(P\) be as in the preceding theorem. If a similarity takes \(P\) into \(\mathbb Q^2\), then \(D(P)\ge\lceil n/2\rceil\).
\end{theorem}

\begin{proof}
After a further scaling the coordinates are integers. Split a node according to the parity of \(x+y\). Modulo two the squared distance is the sum of the two points' parities, so different target classes give different distances from any pin. Within either child, translate an integer point of the child to the origin and apply
\((x,y)\mapsto((x+y)/2,(y-x)/2)\).
The output is integral and squared distances are halved. Apply the transformation also if the split has only one nonempty child. Every surviving pair has a positive integer squared distance, so no branch can halve indefinitely. Suppressing one-child nodes yields the tree required by 982-D01.
\end{proof}

\begin{proposition}[982-D03, 982-D04]
Let \(K\subset\mathbb R\) be a number field. If the quadratic extension \(K(i)/K\) ramifies at some finite place above \(2\), then every \(P\subset K^2\) as in 982-D01 satisfies \(D(P)\ge\lceil n/2\rceil\). In particular this holds for every odd-degree real number field, and the conclusion is invariant under a common similarity.
\end{proposition}

\begin{proof}
At a ramified place the local quadratic extension \(L=K_v(i)/K_v\) is totally ramified. Its residue field has characteristic two, and conjugation \(\sigma:i\mapsto-i\) acts trivially there. Identify \((x,y)\) with \(x+iy\); the squared distance for differences is the norm \(N(z)=z\sigma(z)\).
For each nonsingleton group translate by a member and divide by a power of a uniformizer of \(L\) chosen so that the minimum valuation among all nonzero differences becomes zero. This multiplies all norms of differences by the same nonzero element of \(K_v\). Reduction now gives at least two classes. If the pin has residue \(a\), then the norms of differences to targets with residues \(b\) and \(c\) reduce to \((a-b)^2\) and \((a-c)^2\). These differ for \(b\ne c\). Binary regrouping of distinct residue classes preserves separation. Each original residue child is strictly smaller. This proves the first claim.

At a place where \(v(2)\) is odd, the extension of the normalized valuation satisfies \(v(1+i)=v(2)/2\), since \((1+i)^2=2i\) and \(i\) is a unit. The nonintegral value forces ramification of degree two. In an odd-degree number field \(\sum_{\mathfrak p\mid2} e_{\mathfrak p} f_{\mathfrak p}=[K:\mathbb Q]\), so an odd \(e_{\mathfrak p}\) exists.
\end{proof}

\begin{corollary}[982-D05, 982-D06]
If a strictly convex real counterexample exists at size \(n\), one exists at that size with real algebraic coordinates. Every real number field containing the coordinates of a similar copy of any counterexample has even degree, all its ramification indices above \(2\) are even, and \(K(i)/K\) is unramified at these places.
\end{corollary}

\begin{proof}
For a counterexample, prescribe the signs of all oriented triple determinants and of all within-pin differences of squared distances. The prescriptions are a finite existential formula over \(\mathbb Q\). The transfer theorem for real closed fields gives a solution over the real algebraic numbers, preserving the supporting edges and all within-pin distance partitions. The local criterion and the valuation calculation above give the restrictions on every field containing the coordinates.
\end{proof}

\section{Rational cocircular sets}

\begin{theorem}[982-C06, 982-C07]
Among \(m\ge3\) distinct rational-coordinate points on a circle, the minimum of \(D(P)\) is \(m-1\) for \(m\in\{3,5,6,7\}\), and \(m-2\) for \(m=4\) and \(m\ge8\). At \(m=8\) the sets attaining \(m-2\) are precisely the unions of two disjoint rational squares on the circle.
\end{theorem}

\begin{proof}
The center is rational by solving the two linear bisector equations for three of the points. Translate to the center and divide complex coordinates by one sample coordinate. The normalized coordinates lie in \(\mathbb Q(i)\) and have norm one. The finitely generated multiplicative group they generate has finite torsion contained in \(\mu_4=\{1,i,-1,-i\}\), since the roots of unity in \(\mathbb Q(i)\) are algebraic integers in \(\mathbb Z[i]\) of norm one. Write \(v(z)\) for the image of a sample coordinate in the free quotient, regarded as a vector in a real space.

At a pin \(p\), equal distances to distinct \(b,c\) imply \(bc=p^2\), hence \(v(b)+v(c)=2v(p)\). If \(v(p)\) is an extreme point of the hull of the sample images, both targets belong to its fiber, a subset of \(p\mu_4\). There can be only the one repeated pair \(ip,-ip\) there, so \(D(P)\ge m-2\). If no pin distinguishes all other points, every extreme fiber must be a whole square: once a point is in the fiber both its quarter turns must be present. If the free quotient has positive rank then its generating sample images include zero and have at least two extreme points. Hence \(m\ge8\), and equality at \(m=8\) requires precisely two whole extreme fibers. If the free quotient has rank zero, there are at most four points. A square gives equality at \(m=4\).

For \(m=8+4k+r\), \(0\le r\le3\), let \(w=(3+4i)/5\). It is not a root of unity. Take the two full squares \(w\mu_4,w^{-1}\mu_4\), also the \(k\) squares \(w^j\mu_4\) for \(2\le j\le k+1\), and any \(r\) points of \(\mu_4\). These are disjoint cosets with rational coordinates. Each point in a square has an equidistant pair of neighbors. Each added unit \(z\) sees \(wz,w^{-1}z\) at equal distance, so the lower bound is sharp.
\end{proof}

\section{A sharp annular condition}

\begin{theorem}[982-M06]
Let \(P\) consist of at least three points with distinct polar directions, minimum
angular gap \(\delta>0\), and radii in \([a,b]\) with \(a>0\). If
\[
 b/a<\sqrt{3+\cos^2\delta}-\cos\delta,
\]
then at every pin the distances strictly increase with angular distance on each
half-circle of directions. Each positive distance has multiplicity at most two.
The radius threshold for this multiplicity property is sharp for \(0<\delta\le\pi/2\).
\end{theorem}
\begin{proof}
Let the pin have radius \(t\) and angle zero. For two targets of radii \(r,s\) at angles
\(\theta,\theta+\Delta\) in one half-circle, the difference of squared distances is
\[
 s^2-r^2+2t\{(r-s\cos\Delta)\cos\theta+s\sin\Delta\sin\theta\}.
\]
For \(0<\Delta<\pi\) the expression in braces has a single critical point in
\((0,\pi)\), a maximum, since the sine coefficient is positive. The minimum on
\([0,\pi-\Delta]\) is at an endpoint. For \(\Delta=\pi\) the interval is a singleton.
At zero the difference is at least \(2a^2(1-\cos\delta)-(b-a)^2>0\).
Indeed, with \(z=\sin(\delta/2)\), the positive root of
\(X^2+2X\cos\delta-3\) is smaller than \(1+2z\): at that test value the polynomial
is \(8z(1-z^2)>0\). At the other endpoint the difference is at least
\(a^2-r^2+2t(a-r\cos\delta)\), concave as a function of \(r\).
At \(r=a\) the lower bound is positive. At \(r=b\) we have
\(2(1-(b/a)\cos\delta)>(b/a)^2-1\ge0\), so \(t\ge a\) gives
\(3a^2-b^2-2ab\cos\delta>0\).

For sharpness take unit points at angles \(0,\pi\), and points at angles
\(\pi-\delta,\pi+\delta\) with radius \(b=\sqrt{3+\cos^2\delta}-\cos\delta\).
The minimum gap is \(\delta\) for \(\delta\le\pi/2\). The pin at angle zero
sees all three other points at distance two. The four points are strictly
convex: the outer points have strict radial supports, and the two unit points
do as well because \(b\cos\delta<1\).
\end{proof}

\section{Constant width and unimodal distances}

\begin{theorem}[982-S03, 982-S04]
Let \(P\) consist of \(n\ge3\) distinct boundary points of a planar convex body of constant width \(w>0\). At least \(n/2\) of the points for even \(n\), and at least \(\lceil2n/3\rceil\) for odd \(n\), see \(\lfloor n/2\rfloor\) or more distinctions. Every pin with a unique supporting line has at most two occurrences of each positive distance.
\end{theorem}

\begin{proof}
All point-pair distances in the body are at most \(w\), by projection in the direction of their difference. For a supporting line at \(q\) with interior unit normal \(u\), any point on the opposite supporting line has projection difference \(w\) from \(q\); since the full length is at most \(w\), this point is \(q+wu\). Applying the same argument at the opposite line shows each support face is a singleton. Hence the body is strictly convex.

Fix a pin \(p\). At a local extremum \(q\ne p\) of the boundary distance to \(p\), the one-sided derivatives of \(|q-p|^2/2\) with respect to oriented arc length enclose zero. They are the dot products of \(q-p\) with the one-sided unit tangents. The tangents between these two directions define supporting lines at \(q\). One has zero dot product by continuity. Thus \(p-q\) points along the interior normal to a supporting line, and by the previous paragraph \(p=q+wu\). Consequently there is no local extremum at a positive distance less than \(w\). Starting and ending at \(p\) yields strict increase to the first point at distance \(w\), a possibly trivial plateau, then strict decrease. At a unique-support pin, there can only be one antipodal point at distance \(w\): a chord of length \(w\) attains full width along its own direction and is normal to a supporting line at both ends. This proves the assertion about multiplicity at that pin.

For a finite sample, the boundary list of positive distances at a pin is unimodal. Let \(t_p\) count consecutive equalities in that list. After collapsing constant blocks, each value occurs at most twice and the maximum only once, giving \(n-1-t_p\le2d_P(p)-1\). Every counted equality corresponds to a polygon edge whose bisector goes through the pin. Since the bisector crosses the boundary at the edge midpoint and at most once more, each edge contributes at most one, and \(\sum_p t_p\le n\). An even-size failing pin needs \(t_p\ge2\); an odd-size failing pin needs \(t_p\ge3\). This proves both counts.
\end{proof}

\section{The two endpoints of a unique diameter}

\begin{theorem}[982-X01]
Let \(A,B\) form the unique diameter of a strictly convex \(n\)-gon, \(n\ge3\), with \(s\ge t\) other vertices on the two sides of its line. Then \(d_P(A)+d_P(B)\ge s+3\). For fixed \(s,t\), every pair of integers \(u,v\ge2\) with \(u+v=s+3\) occurs as the two counts. The minimum possible \(\max(d_P(A),d_P(B))\) over all such polygons of size \(n\) is \(\lceil n/4\rceil+1\).
\end{theorem}

\begin{proof}
Normalize \(A=(0,0), B=(D,0)\). Strictness of the diameter forces the other vertices into the open strip \(0<x<D\). On the longer upper chain, let vertex \(j\) be one of maximum height, indexing the \(s\) nonendpoints in increasing order of their \(x\)-coordinates. The chain is concave. The distances from \(A\) increase strictly through vertex \(j\) as both coordinates are nondecreasing on that part and \(x\) increases strictly. From vertex \(j\) to the last, the distances to \(B\) strictly decrease as \(D-x\) strictly decreases and the height is nonincreasing. Adding the diameter gives \(d_P(A)\ge j+1\), \(d_P(B)\ge s-j+2\).

For sharpness take \(A=(0,0), B=(10,0), Y=(1,4)\). Choose \(k\) points on \((x-10)^2+y^2=97\) with \(0.99<x<1,\ y>0\), and \(l\) points on \(x^2+y^2=17\) with \(1<x<1.01,\ y>0\), where \(k+l=s-1\). Use \(Y\) and these points for the upper chain, and reflect any \(t\) of them across \(AB\) for the other chain. Both arcs are strictly concave, with left and right derivatives \(9/4\) and \(-1/4\) at \(Y\). The outer secants have slopes respectively greater than \(3.9\) and less than \(-3.9/9\) whenever the respective arc contributes a point; the arc slopes there lie respectively below \(9.01/3.9\) and above \(-1.01/3.9\). These separate each chain's secant slopes strictly, including a missing arc since the slopes of \(AY,YB\) are \(4,-4/9\). The full reflected polygon is strictly convex, and any desired subset of the reflections inherits that property.

At an upper-chain point left of \(Y\) the squared distances from \(A,B\) are \(20x-3,97\); right of \(Y\) they are \(17,117-20x\). At \(Y\) they are \(17,97\). All are less than \(100\). The squared distance between nonendpoints is at most \(8^2+0.02^2<100\). The diameter is unique, and \(d_P(A)=k+2,\ d_P(B)=l+2\). Choose \(k=u-2,l=v-2\) for arbitrary sharp pairs. Balancing \(u,v\) and minimizing \(s\) at \(\lceil(n-2)/2\rceil\) gives the claimed minimum.
\end{proof}

\section{An interlaced polygon with average below half}

\begin{proposition}[982-X04]
For every \(s\ge1\) there is a strictly convex polygon with \(n=4s+2\)
vertices, half of which see \(n/2-1\) distances and half of which see \(n/2\).
Its total of pinned counts is \(\binom n2\).
\end{proposition}
\begin{proof}
Let \(\alpha=\pi/(2s+1)\) and place points on rays of consecutive angles
\(j\alpha\), using radius \(1\) at even \(j\) and
\(b=2\cos(\alpha/2)-1\) at odd \(j\). Then \(c=\cos\alpha<b<1\).
The cyclic polygon is simple by the angular sectors. The interior point
of each consecutive triple lies strictly outside the opposite chord,
since an inner vertex has radius above \(c\) and an outer vertex has
radius above \(bc\). Hence the polygon is strictly convex.

At an outer pin the within-layer distances have angular bins
\(2\alpha,4\alpha,\ldots,\pi-\alpha\), with distance \(L(\theta)=2\sin(\theta/2)\).
The interlayer bins are \(\theta=\alpha,3\alpha,\ldots,\pi-2\alpha\), and \(\pi\),
with squared distance \(D_b(\theta)^2=1+b^2-2b\cos\theta\).
For bins below \(\pi\) one has
\(L(\theta-\alpha)<D_b(\theta)<L(\theta+\alpha)\).
For the lower inequality first take \(b=c\); the squared difference is
\(\sin\alpha(2\sin\theta-\sin\alpha)>0\), and \(D_b\) increases with \(b\)
here. For the upper inequality \(D_b(\theta)<L(\theta)\) follows by
factoring the squared difference as \((b-1)(b+1-2\cos\theta)\).
At \(\pi\) the distance \(1+b=L(\pi-\alpha)\) duplicates the largest
outer-layer distance. This gives exactly \(2s\) distances.

At an inner pin the within-layer bins have length \(bL\). In each interlayer
bin below \(\pi\) one has \(bL(\theta)<D_b(\theta)<bL(\theta+\alpha)\).
The lower inequality follows from
\(D_b(\theta)^2-b^2L(\theta)^2=(1-b)(1+b-2b\cos\theta)>0\). For the upper
inequality \(D_b(\theta)^2/b^2\) decreases with \(b\); at \(b=c\), the
squared gap before division by \(c^2\) is
\(\sin\alpha(2\cos\alpha\sin(\theta+\alpha)-\sin\alpha)\ge0\), since
\(\alpha\le\pi/3\) and \(\sin(\theta+\alpha)\ge\sin\alpha\) in these bins.
The separate antipodal interlayer distance \(1+b>2b\) contributes the
largest bin. Thus the inner pin sees exactly \(2s+1\) distances. The
two layers have equal size, giving the stated sum.
\end{proof}

\section{Exceptional metric parameters}

\begin{theorem}[982-A02]
Let \(P\subset K^2\) be strictly convex with \(n\ge3\), where \(K\subset\mathbb R\) is a subfield. Let \(Q\) be a real positive definite symmetric matrix with no nonzero scalar multiple over \(K\). In the \(Q\)-metric, every pin sees multiplicity at most two. Some pin sees at least \(n-3\) positive distances, and at least \(n-2\) if \(n\) is odd.
\end{theorem}

\begin{proof}
For distinct \(p,q,r\) put \(u=q-p\), \(v=r-p\). Equal distances at \(p\) imply \(\operatorname{tr}(QB)=0\) for \(B=uu^{\mathsf T}-vv^{\mathsf T}\in\operatorname{Sym}_2(K)\). Because the three entries of \(Q\) span at least two dimensions over \(K\), the kernel over \(K\) of this linear map has dimension at most one. Thus all actual nonzero equality matrices are proportional. Three equidistant targets at one pin would give two independent matrices: in the basis of the first two target vectors the third has nonzero coordinates along both axes, so its rank-one matrix has nonzero off-diagonal entry. This proves the multiplicity bound.

If any equality exists, its matrix defines the product of the two linear forms given by \(u-v\) and \(u+v\). They are nonproportional since no three vertices are collinear. Uniqueness of the factorization fixes an unordered pair of directions common to all equalities. Each possible base \(qr\) lies along one of the two directions. At most \(\lfloor n/2\rfloor\) vertex pairs have one fixed direction, and each pair has at most two \(Q\)-bisector witnesses. Hence there are at most \(4\lfloor n/2\rfloor\) equal-target pairs counted over all pins. Averaging gives a pin with at most two such pairs for even \(n\) and at most one for odd \(n\). Since no class has more than two targets this gives the stated distance counts.
\end{proof}

\begin{corollary}[982-A06]
For a fixed strictly convex \(n\)-point set, up to scale there are at most
\(\binom{n-1}{3}\) positive definite metrics that need testing for a counterexample.
\end{corollary}
\begin{proof}
Fix one pin. If every distance class there has size at most two, that pin reaches
the half bound. Otherwise some triple of targets is equidistant. Its equality
matrices are independent by the preceding basis argument. Hence its two
linear equations determine at most one metric up to scale. Enumerate the
triples and retain only projectively positive definite forms.
\end{proof}

\begin{theorem}[982-A01]
For a fixed strictly convex planar set of \(n\ge3\) points there are at most \(O(n)\) positive definite symmetric forms modulo positive scaling for which \(\sum_p d_P(p)<\binom n2\). The constant is absolute.
\end{theorem}

\begin{proof}
Normalize \(\operatorname{tr}Q=1\) so that the forms live in an affine plane. Each target pair with a distinct pin defines the event \(\operatorname{tr}(QB)=0\). In this plane it is a line: \(B\) is nonsingular indefinite, since its two target vectors are independent. Discard lines containing no positive definite form. Coincident lines have proportional matrices because the corresponding homogeneous planes are determined by the normalized line and the origin. The factorization above constrains each base direction to a fixed pair of directions on that line. At any one positive definite form on that line, at most \(2\lfloor n/2\rfloor\) bases occur and each has at most two qualifying pins. This gives weight at most \(w=2n\) per line and total weight at most \(W=n\binom{n-1}{2}\).

Let \(X\) be the cardinality of any finite subset of exceptional forms and \(K_0=\binom n2\). Over all pins every exceptional form loses more than \(K_0\) counts to repeated distances. Each lost count can be charged to a different equal-target pair at its pin. Thus total weighted incidences exceed \(K_0X\). Apply the Szemer\'edi--Trotter bound \cite{st} to the unweighted set of lines of weight at least \(t\), and integrate for \(0<t\le w\). The number of those lines is at most \(W/t\), its integral is \(W\), and the integral of its two-thirds power is at most \(3 W^{2/3}w^{1/3}\). This gives
\[
 K_0X\le C_0\bigl(X^{2/3}W^{2/3}w^{1/3}+Xw+W\bigr)
\]
with an absolute \(C_0\). For large enough \(n\) absorb \(C_0Xw\); either the first or last term then bounds \(X\) by \(O(wW^2/K_0^3+W/K_0)=O(n)\).

For the finitely many smaller \(n\), a form violating the sum bound has a class of size at least three at some pin. Otherwise \(d_P(p)\ge\lceil(n-1)/2\rceil\) at every pin. Each such triple of targets imposes two independent equality matrices, by the same basis argument as above. This determines the form up to scaling. There are at most \(n\binom{n-1}{3}\) possibilities, a uniform multiple of \(n\) over the bounded range. Since the bound applies to any finite selection of exceptional forms, it bounds the total. If every pin sees fewer than \(\lfloor n/2\rfloor\), the sum is below \(K_0\), so the same conclusion holds for the half-bound violations.
\end{proof}

\section{Multiplicity identities}

\begin{proposition}[982-X07, 982-X08]
Let \(s\) be the number of singleton distances at \(p\), and \(e\) the sum of the excesses \(k-2\) over all distance classes at \(p\) with \(k\ge3\). Then \(2d_P(p)=n-1+s-e\) and the pin reaches \(\lfloor n/2\rfloor\) exactly when \(s\ge e\). If a strictly convex \(P\) is a counterexample and \(e_j\) counts its targets pairs with \(j\) bisector witnesses, then \(e_2-e_0\ge3n/2\) for even \(n\) and \(e_2-e_0\ge3n\) for odd \(n\).
\end{proposition}

\begin{proof}
Summing multiplicities gives \(n-1=2d_P(p)-s+e\). If \(n\) is even, the threshold initially reads \(s-e\ge1\), equivalent to \(s-e\ge0\) because \(s-e\) is odd in this case. For odd \(n\) the threshold reads \(s-e\ge0\) directly.
Each pair has zero, one, or two witnesses, so the apex-counted isosceles total is \(\binom n2+e_2-e_0\). Since \(\binom k2\ge2k-3\), it is at least \(2n(n-1)-3\sum_p d_P(p)\). Substitute \(\sum_p d_P(p)\le n(\lfloor n/2\rfloor-1)\) and simplify in the two parities.
\end{proof}

\section{Points covered by polynomial graphs}

\begin{theorem}[982-P03]
Let \(P\) be \(n\ge2\) distinct points covered by \(r\ge1\) polynomial graphs in a common orthonormal coordinate system, all of degree between \(2\) and \(d\). There is a constant \(C_d\) depending only on \(d\) such that \(D(P)\ge n-1-C_d r^{9/7}n^{5/7}\). For fixed \(r,d,\epsilon>0\) the number of pins missing at least \(\epsilon n\) distances is \(O_{r,d,\epsilon}(\sqrt n)\).
\end{theorem}

\begin{proof}
For graphs \(f,g,h\) of degrees \(a,b,c\) and leading coefficients \(A,B,C\), define \(D_g(x,y)=(y-x)^2+(g(y)-f(x))^2\) and \(E(x,y,z)=D_g(x,y)-D_h(x,z)\). Over \(\mathbb C\) the coefficients of \(E\) as a polynomial in \(x\) have common divisor at most \(y-z\), by looking at degree \(a\) then degree one. This factor occurs exactly once when \(g=h\), and is absent otherwise. Remove it when present. The resulting \(F\) still vanishes on all relevant triples of distinct targets.

To prove absolute irreducibility, let \(L\) be the least common multiple of \(a,b,c\). Substitute \(y=t^{-L/b},z=v t^{-L/c},x=w t^{-L/a}\) into \(E\), and multiply by \(t^{2L}\). Its specialization at \(t=0\) is \((B-Cv^c)(B+Cv^c-2Aw^a)\). Over \(\mathbb C(v)\) the first factor is a unit; the second is Eisenstein in \(w\) for the valuation of \(v\) at any simple zero of \(B+Cv^c\). Both the original polynomial in \(w\) and its specialization have degree \(a\), with unit leading coefficient in the valuation ring at \(t=0\). A factorization into positive \(w\)-degrees would therefore specialize to one with the same degrees. The substitution on \(\mathbb C(y,z)\) is injective, so Gauss's lemma and the preceding primitive-part computation give irreducibility of \(F\).

For large positive \(x\), the critical target equation
\[
 y-x+(g(y)-f(x))g'(y)=0
\]
has a complex analytic branch \(y_*(x)\sim\lambda x^{a/b}\), where \(B\lambda^b=A\). To see this, replace \(x=t^{-b},y=t^{-a}Y\) and multiply by \(t^{a(2b-1)}\). The constant term in \(t\) is \((BY^b-A)bBY^{b-1}\), with simple roots at such \(\lambda\). The implicit function theorem supplies analytic \(Y(t)\). The convergent expansion shows \(y_*'\ne0\) for all sufficiently large \(x\). The critical equation gives \(g(y_*)-f(x)=(x-y_*)/g'(y_*)\) and hence \(S(x):=D_g(x,y_*)=o(|f(x)|^2)\). Also \(D_{g,yy}/2=1+g'^2+(g-f)g''\sim g'^2\ne0\) at this branch.

In the equation \(D_h(x,z)=S(x)\), all complex roots satisfy
\[
 |h(z)-f(x)|\le\sqrt{|S(x)|}+|z|+|x|.
\]
Leading-term growth forces \(|z|=O(|f(x)|^{1/c})\), uniformly over roots for the fixed graphs. Then \(h(z)/f(x)\to1\) and \(|z|\asymp|f(x)|^{1/c}\). Thus \(D_{h,zz}/2\sim h'(z)^2\ne0\), and each root has multiplicity at most two. If all roots were double, \(D_h(x,z)-S(x)=W(z)^2\) with \(W\) chosen to have the leading coefficient of \(h\). Then
\[
 [W-(h-f(x))][W+(h-f(x))]=(z-x)^2-S(x).
\]
The second factor has degree \(c\). This excludes \(c>2\). If \(h(z)=Cz^2+C_1z+C_0\), the first factor must be a nonzero constant \(u\), forcing \(2uC=1\) and \(2uC_1=-2x\), impossible for all sufficiently large \(x\). A simple root therefore exists at every sufficiently large \(x\). A local simple branch \(z_*(x)\) cannot be constant on an interval. Otherwise analytic continuation in the real parameter along the critical branch gives \(D_h(x,z_0)=S(x)\) up to infinity, but the left side divided by \(f(x)^2\) tends to one.

Choose a point with \(z_*'\ne0\). There \(E_z\ne0,E_y=0,E_x=-E_z z_*'\ne0\). If \(g=h\), the simple root is distinct from the critical target. In a neighborhood write \(z=H(x,y)\); \(H_y/H_x=E_y/E_x=(y-y_*(x))U(x,y)\) with \(U\) nonzero and analytic, since \(E_y\) is independent of \(z\) and \(E_{yy}\ne0\). The mixed logarithmic derivative \(\partial_x\partial_y\log(H_y/H_x)\) has a double pole with coefficient \(y_*'(x)\). In an additive local representation it would vanish, since the ratio would separate by variables. The derivative is rational on the irreducible surface, so a generic additive representation would force its vanishing on the whole surface where defined. This contradiction, and generic nonvanishing of the three partials, eliminate the additive and cylindrical exceptions of Solymosi--Zahl \cite[Theorem 1.4]{sz}.

Assign each point to exactly one graph and write \(n_i\) for the sizes of the nonempty groups. For each ordered triple of group roles, let \(u\le v\le w\) be their sizes. The incidence bound is \(O_d((uvw)^{4/7}+v\sqrt w)\). The first term sums to at most \((\sum_i n_i^{4/7})^3\le r^{9/7}n^{12/7}\). The second sums to at most \(6r(\sum_i n_i)(\sum_i\sqrt{n_i})\le6r^{3/2}n^{3/2}\), absorbed by the first bound using \(r\le n\). The total loss in distance counts is bounded by the number of equal-target pairs. Averaging over pins proves the first conclusion. If only the \(s\) pins with loss at least \(\epsilon n\) are used, for fixed \(r\) the same incidence bound gives \(\epsilon ns\ll_{r,d}(sn^2)^{4/7}+n^{3/2}\), proving the exceptional-pin estimate.
\end{proof}

\section{A strict majority on a circle}

\begin{theorem}[982-C11 and the strict convex part of 982-C10]
If \(P\) is a strictly convex planar counterexample and \(m\) of its \(n\) points lie
on a circle, then \(m\le n/2+O(n/\log(2n))\), with an absolute constant.
At every size \(m\ge 2(n-m)+42\) forces a qualifying pin on the circle.
\end{theorem}

\begin{proof}
Set \(k=n-m,\ \Delta=m-k>0\), as there is nothing to prove for the asymptotic statement if \(m\le n/2\).
Normalize the circle to the unit circle and one vertex of its set \(A\) to \(1\).
If \(A\) is contained in a closed semicircle, an end vertex sees \(m-1\) distances along it, reaching the target. Otherwise the center is inside \(\operatorname{conv} A\), and cannot be another vertex.
The group generated by \(A\) has free rank \(r\). The images of \(A\) modulo torsion affinely span \(\mathbb R^r\). For a pin in an extreme fiber of size \(t\), the relation \(bc=a^2\) for equal chord lengths forces both targets of each coincidence to be in the same fiber. The pin therefore sees at least \(m-\lceil t/2\rceil\) radii. If only the circle pins are assumed to fail, every extreme fiber has \(t\ge\Delta+1\), and
\[
 (r+1)(\Delta+1)\le m. \tag{A}
\]
For the asymptotic claim all pins fail. Each off-circle pin \(q\) has a repeated radius among its \(m\) circular targets, so \((q/|q|)^2=bc\) for some \(b,c\in A\). All directions belong to the group \(\Gamma\) of square roots of the group generated by \(A\), still of rank \(r\).

\paragraph{Uniform equalities at different radii.}
Apply Evertse--Schlickewei--Schmidt \cite{ess} to sums of fixed numbers of Laurent monomials with variables in \(\Gamma^d\). Divide each minimal vanishing subsum by one of its monomials. Each resulting nondegenerate equation in ratios has at most \(\exp(O(r+1))\) solutions, since the rank of the monomial image is at most \(dr\). Fixing those ratios cuts out a translate of a finite-component subtorus. Its number of components depends only on the exponent matrix by Smith normal form. Each component with a group point has such a point as a representative. Thus the solutions are covered by \(\exp(O(r+1))\) translates on which the equation is identically zero.
Put \(M=\exp(C(r+1))\), for a fixed \(C\) large enough for the few exponent patterns below. The equation
\[
 \alpha U+\beta U^{-1}+\gamma V+\delta V^{-1}+e=0
\]
with all five coefficients nonzero has no positive-dimensional component of this kind: if both coordinates vary, the constant cannot cancel; if just one varies, its distinct character and inverse cannot cancel. Thus there are at most \(M\) solutions. Squaring distances gives this equation for a fixed noncentral pin and an unknown pair of targets at two fixed unequal radii. It also gives this equation for a fixed target \(q\ne0\), a variable unit-circle pin \(z\) and an unknown target \(Tvz\) on a radius \(T\ne|q|\), using variables \(z,v\). In the latter case at most \(M\) circle pins create any such coincidence.

\paragraph{Double witnesses.}
Fix an off-circle target \(q\) and rotate so \(q=R>0\). If another off-circle target \(b\) has the same distance as \(q\) at two different circle pins \(z,w\), reflection in their chord gives \(b=z+w-Rzw\). With \(\lambda=b/|b|\) set \(X=z/\lambda,\ Y=w/\lambda,\ Z=zw/\lambda\). The reality of \(b/\lambda\) gives
\[
 X-X^{-1}+Y-Y^{-1}=R(Z-Z^{-1}). \tag{B}
\]
On a torus translate where (B) holds, classify the nonconstant coordinate characters up to inversion. Any nonzero class must have at least two coordinates in it. Thus at most one class occurs. A constant triple fixes both pins. If only \(X,Y\) vary, then \(Z=\pm1\) and either \(X+Y=0\) or \(XY=1\); these give respectively equal target radii or collinearity with the center. If only \(X,Z\) or \(Y,Z\) vary, cancellation forces \(R=1\), excluded. If all three vary, write \(Z=st,\ X=xt^{\epsilon_1},\ Y=yt^{\epsilon_2}\) with \(\epsilon_i=\pm1\) and unit modulus coefficients. A positive exponent fixes \(w=Z/X\) or \(z=Z/Y\). When both are negative, comparison of the coefficient of \(t\) gives \(x^{-1}+y^{-1}=-Rs\) and \(xy=s^{-2}\) by conjugating. Substitution gives \(b/q=\sigma-\sigma^2\) with \(\sigma=-s^2t^2\) of unit modulus.

Collect the fixed pins in an exception set \(E_q\) of size at most \(M\), increasing the absolute constant if necessary. Repeating the argument with \(q,b\) interchanged, outside the equal-radius, collinear and exceptional cases both \(v=b/q\) and \(1/v\) would have the last form. Writing \(H=|v|^2\) gives \(2\Re v=3H-H^2=3-1/H\), so \((H-1)^3=0\) and \(b=q\), a contradiction.

\paragraph{Layer counting.}
For \(s\) off-circle targets with \(s_j\) on each radius layer, count losses at circle pins by connecting each distance class of those targets with a simple path. Each unordered pair appears at most twice over all pins. Its first occurrence costs at most \(\binom s2\) in aggregate. Second occurrences in the same radius cost at most \(\sum_j\binom{s_j}{2}\); centered collinear pairs cost at most \(s/2\). Other second occurrences have an exceptional pin on their path edge, and at each \((q,z)\) at most two edges meet \(q\). Thus total loss is at most \(\binom s2+\sum_j\binom{s_j}{2}+(2M+1)s\).

Call a layer large if it contains at least \(\Delta/4\) points; there are \(L\le4n/\Delta\) large layers and \(s\) small-layer points. Count radii including zero. The large layers have at least \((n-s)/2\) distances per pin before merging, as every within-layer multiplicity is at most two. Cross-layer losses at a fixed pin cost at most \(M\binom L2\). Each small-layer target can merge with a large layer at no more than \(M\) circle pins. Since each circle pin sees at most \(n/2\) distances including zero,
\[
 ms/2\le\binom s2+\sum_j\binom{s_j}{2}+(2M+1)s+mM\binom L2+sML .
\]
Here \(m-s\ge\Delta\) and \(s_j<\Delta/4\). The left after moving the first two terms is at least \(3\Delta s/8\). From (A), \(M\), \(s\), and \(L\) are each bounded by \(\exp(O(n/\Delta))\).

For every pin, reflection in its radial axis misses at most \(E=s+2M\binom L2\) targets in any one large layer: every missed point produces a singleton distance before merging the large layers, and the total count is at most \(n/2\). Compose each reflection with that of one fixed pin. There are at least \(n/2\) different rotations since a line through the center contains at most two vertices. On a large layer of size \(v\), each rotation preserves at least \(v-2E\) points. Distinct rotations cannot send the same ordered pair to one another, so \((n/2)(v-2E)\le v^2\). Either \(v\le4E\) or \(v\ge n/2-4E\). If \(\Delta\le8E\), then \(n=(n/\Delta)\Delta\le\exp(O(n/\Delta))\). Otherwise every large layer other than \(A\) falls into the first option, since all off-circle targets together number \(k=(n-\Delta)/2<n/2-4E\). Thus \(k\le s+4EL\le\exp(O(n/\Delta))\).

At a circle pin, start with at least \(\lfloor m/2\rfloor\) circle distances then adjoin the \(k\) targets. At least \((k+1)/2\) counts must merge if the pin fails. Off-circle targets merging with circle distances cost at most \(Mk\) in aggregate by the five-term equation with \(T=1\); off-circle pairs cost at most \(2\binom k2\). This gives \(m(k+1)/2\le Mk+k(k-1)\), hence \(m\le2M+2k\). In all cases \(n\le\exp(O(n/\Delta))\), proving the asymptotic estimate.

For the explicit threshold alone assume just that the circle pins fail, so (A) still applies and \(m\ge2k+42\) forces \(r=0\). We can bound the cost of adding an off-circle target against existing circle distances directly. With \(u=a/z\) for a circle target \(a\) and pin \(z\), the equation is
\[
 u+u^{-1}-\bar q z-qz^{-1}+|q|^2-1=0.
\]
This has finitely many torsion solutions by the same five-term argument applied to the rank-zero group of all roots of unity. Its Newton polygon has area two, so Beukers--Smyth \cite{bs} bounds the number by 44. The solution \(u=1\) is impossible; \(u=-1\) occurs at at most two pins by intersection of the sample circle with the radius-two circle centered at \(q\). The other solutions occur in pairs \(u,u^{-1}\), so at most 23 pins have a merge with this target into the circle distances.
Now \(m(k+1)/2\le23k+k(k-1)=k(k+22)\), contradicted by \(m\ge2k+42\).
\end{proof}

\section{Strictly convex norms}

\begin{theorem}[982-T02]
For \(n\ge4\) planar points in strictly convex position and any strictly convex norm, some pin sees at least \(\lfloor n/3\rfloor+1\) norm distances.
\end{theorem}
\begin{proof}
The strict quadrilateral crossing inequalities hold under the norm, by splitting intersecting diagonals and applying the strict triangle inequality. Thus each oriented pair \(ij\) has at most one equidistant witness on its forward open chain. Indeed two such witnesses would make the sum of the diagonals equal to the sum of the corresponding opposite sides. Let \(\chi(i,j)\) be the number of witnesses on that chain.

We use a normalized angular measure on the norm unit circle. Two distinct circles in a strictly convex normed plane meet in at most two points \cite[Proposition 14]{msw}. For each unit \(u\) let \(T(u)\) be the unique unit \(v\) on the positively oriented open semicircle with \(\|v-u\|=1\). Existence follows by continuity on this semicircle, and uniqueness because \(u-v\) also lies on both circles and on the other semicircle. The successive directions are \(u,v,v-u,-u,-v,u-v\); in particular \(T^3u=-u\) and \(T^6u=u\). Compactness and uniqueness give continuity, and the periodic identities show it is an orientation-preserving homeomorphism. Average ordinary angular measure over its six powers. This measure is nonatomic, strictly positive on open arcs, and assigns mass \(\pi/3\) to each arc from \(u\) to \(T(u)\). At a consecutive equilateral triple of the polygon, the external turn at its middle point has mass \(2\pi/3\). With \(n\ge4\) there are at most two such turns, since the external turns partition total mass \(2\pi\) and all have positive mass.

Besides the \(n\) missing witnesses for consecutive forward edges, at least \(\lceil(n-4)/2\rceil\) additional witnesses are missing when \(n\ge5\). In fact write \(x_i=\chi(i,i+2)\) and \(y_i=\chi(i,i+3)\) with cyclic indices. If \(x_i=x_{i+1}=y_i=1\), then the three consecutive boundary edges at these positions have the same length, and the \(y_i\)-witness is one of the two interior vertices. Its equality forces either the first two or the last two edges to be part of a consecutive equilateral triple. Each middle point accounts for at most two such indices. At most four indices have all three terms equal to one. A zero term is counted at most twice in the remaining indices, proving the missing-witness bound. Thus the total isosceles count is at most \(n(n-2)-\lceil(n-4)/2\rceil\).

If each pin had at most \(k=\lfloor n/3\rfloor\) distances, splitting \(n-1\) targets into balanced classes minimizes the count of equal-target pairs. This minimum is \(n-2,n-1,n+1\) for \(n\equiv0,1,2\pmod 3\) respectively. Summation contradicts the missing-witness bound for \(n\ge5\). For \(n=4\), a single distance at each pin would make all six lengths equal, violating the crossing inequalities.
\end{proof}

\begin{corollary}[982-T03]
Among points on a circle of a strictly convex planar norm every pin sees at least \(\lfloor n/2\rfloor\) distinct norm distances.
\end{corollary}
\begin{proof}
Each positive distance corresponds to intersecting the sample circle with a different norm circle about the pin. The two-circle intersection bound gives at most two targets for each distance.
\end{proof}

\begin{proposition}[982-T01]
For \(3\le n\le9\) the strictly convex norm version of the half bound holds.
\end{proposition}
\begin{proof}
For \(n=3,4,5,6,7,9\) the elementary bound above or the case \(n=3\) suffices. At \(n=8\), if all pins see at most three distances, a class of size \(s\) at a pin forces \(s-1\) different distances at its first target along the boundary: the crossing inequalities between the pin and any three equidistant targets imply strict monotonicity of successive distances from that first target. Thus classes have size at most four.

The archived integer enumeration branches at each pin over every partition of the unprocessed target classes compatible with at most three classes and size at most four. The symmetric undirected edge variables are globally merged by the chosen equalities. A crossing inequality \(X+Y>U+V\) and an existing order \(X\le U\) force \(Y>V\), and likewise with any exchanged left or right term. Neither these strict comparisons nor transitive closure can rule out an actual realization. Every realized partition extends a branch: the initial vertex may be chosen to maximize its descending class-size signature, and reversal canonically chooses one orientation with that signature. The enumerator rejects all branches, with 4,376 nodes in the eight-vertex case. This is an exhaustive combinatorial computation in the supplied Gemini source, with the input and build described in the audit.
\end{proof}

\section{Jacobian rank and local rigidity}

\begin{theorem}[982-J01, 982-J02, 982-J03]
Let \(P\) be \(n\ge3\) distinct planar points with no three collinear. Let \(s\)
count pins with at least \(\lfloor n/2\rfloor\) distances, and let \(J\) be the real
Jacobian of all true equations \(|q-p|^2-|r-p|^2=0\) with distinct \(p,q,r\), evaluated
at \(P\). Write \(t=\dim_{\mathbb R}\ker J\). Then
\[
 \operatorname{trdeg}_{\mathbb Q}\mathbb Q(P)\le t\le 2s+576H_{n-s}.
\]
Here \(H_j=\sum_{i=1}^j i^{-1}\) and \(H_0=0\). By fixing \(t\) suitable scalar
coordinates to their actual values, the remaining coordinates are locally
unique. Over the field \(F\) generated by the fixed values,
\([F(P):F]\le 2^{2n-t}\).
\end{theorem}

\begin{proof}
Each row for a pin \(p\) and targets \(q,r\) has blocks
\[
 2(r-q),\qquad 2(q-p),\qquad 2(p-r),
\]
respectively. At each pin, order each class of targets and retain only
equalities between consecutive targets in that class. Two point blocks then
co-occur in at most six rows. There are at most two adjacencies when one is the
pin and the other a target, for either pin choice. When both are targets their
bisector contains at most two pins, and the target pair is adjacent at most
once at each of them. A fixed target's row block is collinear with its vector
to the pin. Each such direction occurs at most twice among its target roles:
one pin lies on that line and contributes at most two adjacencies.

The Dvir--Garg--Oliveira--Solymosi block rank theorem \cite{dgos} gives
corank at most \(2m/(1+k/24)\) for a matrix on \(m\) two-coordinate blocks
with at most three nonzero blocks in a row, overlap at most six rows between
two columns, and \(k\) well-spread \(1\times2\) blocks in every column. Here
well-spread means that no direction contributes more than \(k/2\), since the
only nontrivial test subspaces in dimension two are lines. All blocks are
real and the matrix rank is unchanged on passing to \(\mathbb C\).

In the velocity space \(\ker J\), freeze all \(s\) qualifying pins at a cost
of at most \(2s\) dimensions. A point equidistant to three frozen points has
velocity zero in the remaining space, since the two independent differences
of those point vectors are orthogonal to its velocity by the actual rows of
\(J\). Freeze such points for free as long as possible. If \(b\) points are
still unfrozen, every unfrozen pin sees at least \((n-b)/2\) distances into
the frozen targets. In the path ordering put all frozen targets first and
then retain only edges whose latter target is unfrozen. Orderings may be
chosen anew at this stage. Each unfrozen pin has at least
\[
 b-1-\bigl(d_P(p)-d_{\rm frozen}(p)\bigr)\ge b/2
\]
retained rows. This follows because every exclusively unfrozen class loses
only its first target from the count of latter targets.

For \(b\ge48\), put \(k=\lfloor b/12\rfloor\). Repeatedly remove an unfrozen
point currently occurring as a target fewer than \(k\) times, together with
its incident rows in any role. This cannot reach \(u=\lfloor b/3\rfloor\)
remaining points. The \(u\) remaining pins started with at least \(ub/2\)
rows and lost at most \((b-u)k\) to deletions: their centers were never
deleted. Their surviving rows number at most \(u(u-1)\) because the latter
target is a different remaining point and is unique at a pin. But for
\(b=3u+r\), \(0\le r\le2\), one has
\[
 12\bigl(u(b/2-u+1)-(b-u)b/12\bigr)=ru+12u-r^2>0.
\]
The surviving \(m>b/3\) points therefore each have \(k\ge4\) target-role
rows, with direction multiplicity at most two. Their velocity projection
lies in the kernel of the corresponding matrix on these \(m\) points,
since the other columns in the retained rows are frozen. Freezing this
batch costs at most \(2m/(1+k/24)\le576m/b\) dimensions. Summing these costs
using \(m/b\le H_b-H_{b-m}\), and charging the final \(b<48\) points directly
at cost at most \(2b\le576H_b\), proves the bound for \(t\).

Derivations with respect to a transcendence basis chosen from the actual
coordinates extend to the finitely generated coordinate field, since the
remaining extension is algebraic and characteristic zero. Applying them to
the equality equations yields independent tangent vectors: their entries
on basis coordinates are the identity. This proves the first inequality
without a smoothness assumption on the entire solution set.

Choose \(t\) coordinate functionals that give an injection on \(\ker J\).
The other \(2n-t\) columns of \(J\) are independent. A square submatrix on
these columns has nonzero determinant, so the corresponding equations after
fixing the \(t\) coordinates define an isolated simple root at \(P\).
The inverse function theorem gives local uniqueness. The derivation argument
over \(F\), applied to this invertible matrix, gives transcendence degree
zero over \(F\). The remaining coordinates are a finite separable extension.
Each embedding over \(F\) produces a different isolated simple complex
solution of the selected equations, all of degree at most two. The affine
B\'ezout bound for isolated solutions gives at most \(2^{2n-t}\) such
embeddings \cite{heintz}.
\end{proof}

\section{Two-circle coverage}

\begin{theorem}[982-M01, 982-M02]
For all sufficiently large \(n\), a strictly convex \(n\)-point set with at least
\(0.99n\) points on the union of two circles has at least \(n/10\) pins attaining
\(\lfloor n/2\rfloor\). If the set is partitioned between two nonconcentric circles
with at least \(n/4\) points in each layer, then at least \(n/2000\) pins exceed
\((1/2+10^{-4})n\). All size thresholds are uniform in the circles.
\end{theorem}

\begin{proof}
Assign each sampled point on the circles to exactly one layer. The proof
uses equalities with one target in each of these disjoint layers.

\paragraph{Mixed equalities.}
For different circles with centers \(d,0\) and radii \(R,S\), use complex
coordinates \(p=d+x,q=d+y,r=z\) for a pin and target on the first circle and
a target on the second. Substituting \(\bar x=R^2/x,\bar y=R^2/y,\bar z=S^2/z\)
in the equality of squared distances and clearing denominators gives
\[
 f=R^2z(x-y)^2+yC,\qquad C=(d+x-z)((dx+R^2)z-S^2x).
\]
As a polynomial in \(y\) it is primitive, since its leading coefficient
is \(R^2z\) and \(C(x,0)=-S^2x(d+x)\ne0\). Its discriminant \(C(C-4R^2xz)\)
has odd valuation at \(d+x-z=0\). On that line the other factor of \(C\)
is \(B(x)=dx^2+(d^2+R^2-S^2)x+dR^2\), which vanishes identically only
for identical circles. Thus \(f\) is irreducible. On the complex arc
\(y=x,z=d+x\), generically \(f_z=-xB(x)\ne0\), \(f_y=0\), \(f_{yy}=2R^2(d+x)\ne0\).
The local graph \(z=g(x,y)\) has \(g_x=1,g_y=0,g_{yy}\ne0\) along \(y=x\).
Hence the rational mixed logarithmic derivative of \(g_x/g_y\) has a
second-order pole there. As in the polynomial-graph proof this excludes
an additive separation. All three variables occur. Separate nondegenerate
Cayley parameter changes reduce the circle parameters to real variables,
preserving irreducibility after removing boundary factors, with bounded
degree. Solymosi--Zahl therefore gives \(O(n^{12/7})\) mixed equal-target
pairs \(J_p\) summed over the circle pins in either layer.

\paragraph{From almost-invariance to a group.}
Call surviving pins good for this step if they fail the desired distance
threshold, lie in a circle layer, are not either center, and satisfy
\(J_p\le n/10000\). There are \(o(n)\) discarded pins due to mixed
equalities and at most two centers. If few pins attain the threshold,
the good set has size at least \(0.889n\) in the coverage problem and
\(0.999n\) in the nonconcentric slack problem.
For each layer of size \(m\), reflection about its center--pin axis has
deficit equal to the number of its targets whose reflected partner is
missing. From the formula for paired distances including zero, the sum
of deficits of the two layers is at most
\(2(d_P(p)+1)-(m_1+m_2)+2J_p\). Thus each deficit is at most
\(\eta=0.011n\) for coverage and \(\eta=0.001n\) for slack.

Fix a good pin \(p_0\) and let \(T\) be the rotations of the layer obtained
by composing each good reflection with its reflection. Then
\(|T|\ge |G|/2\), where \(G\) is the good set, because no line contains
more than two vertices. Each member of \(T\) has deficit at most \(2\eta\);
each of \(H=TT^{-1}\) has deficit at most \(4\eta\). Counting ordered pairs
in the layer gives \(|T|\le m^2/(m-2\eta)\) and \(|H|\le m^2/(m-4\eta)\)
when the denominators are positive.

For coverage, assume each layer has at least \(0.07n\) points. The first
upper bound is at most \(m+3\eta\) since \(m\ge6\eta\), giving
\(0.4115n\le m_i\le0.5885n\). The second upper bound is then at most
\(m+5\eta\le0.6435n\) since \(m\ge20\eta\), and \(0.6435n<3|T|/2\).
For slack, the two layer sizes start at \(n/4\), so the same estimates
give \(0.4965n\le m_i\le0.5035n\) and \(|H|\le0.5085n<3|T|/2\).

In either case \(H\) is a finite subgroup. Indeed \(|T\cap hT|>|T|/2\)
for each \(h\in H\), by translating two copies of \(T\) inside the
difference set. For \(h,k\in H\), a common element of the two more-than-half
subsets implies \(hk\in H\). Thus each layer has a group of order
\(h\ge |G|/2\). If an orbit contains \(b_j\) targets, then
\(\sum b_j^2=\sum_{u\in H}|S\cap uS|\ge h(m-4\eta)\), where \(S\) denotes
the layer directions. Some orbit contains at least \(h(1-4\eta/m)\) points,
exceeding \(0.89h\) for coverage or \(0.98h\) for slack. Every good pin's
squared direction about either center belongs to one fixed phase times
the corresponding group by construction of \(T\).

\paragraph{Distinct centers.}
Normalize the first center to zero and the other to \(d>0\). A first-layer
good pin \(p=Rz\), \(|z|=1\), has second squared direction
\(w=z(Rz-d)/(R-dz)\). Both \(z,w\) are fixed phases times roots of unity.
If \(R\ne d\) the graph has a zero and pole at distinct nonzero finite
arguments. This rules out every nonvertical torsion-binomial component
\(z^aw^b=\omega\) with \(b\ne0\). A vertical component is impossible in
the graph, whose numerator and denominator are coprime. The Newton polygon
of its defining Laurent polynomial has area one, so Beukers--Smyth allows
at most 22 such pins. Each layer actually has a positive linear number
of good pins in both parameter regimes. Applying the argument in both
directions forces each circle's radius to equal \(d\). Their nearly full
orbits ensure each open arc of angular length \(\pi/3\) contains a sample
for large \(n\): there are at least \(h/6-1\) orbit positions and fewer
than \(0.11h\) omissions. After scaling the radii to one, the convex hull
of the second layer contains the disk of radius \(\cos(\pi/6)\) at its
center, by support in every direction. The first circle passes through
that center, and its arc centered there of angular length \(\pi/3\)
contains a first-layer point at distance at most \(2\sin(\pi/12)<\cos(\pi/6)\).
This violates strict convex position and completes the slack result.

\paragraph{Concentric coverage.}
For one circle with \(m\) points, at least \(m-2n/3-14\) pins qualify.
For a semicircle sample, count pins with at least \(\lfloor n/2\rfloor\)
targets farther along one direction. Otherwise the center lies inside
the convex hull. Use the torsion-free projection of the circle direction
group from the majority proof. While the residual circle set of size \(u\)
has at least two extreme fibers and \(u\ge2n/3\), delete the smaller of two
extreme fibers. All its pins qualify since they see at least
\(u-\lceil u/4\rceil\ge\lfloor n/2\rfloor\). If the residual size falls below
\(2n/3\), enough qualifying pins have been counted. Otherwise stop at one
fiber, rotate its coordinates to roots of unity, and write \(s=m-u\) for
deleted qualifying pins and \(t=n-u\). When adding \(t\) targets, those
on the same circle in previously deleted fibers cannot merge with the
original distances, by \(bc=a^2\). Each different-radius target merges with
original distances at at most 23 pins by the torsion count in the majority
proof, and each pair of added targets accounts for at most two merges.
Thus at most \(2t(t+22)/(t+1)\le2t+42\) of the residual circle pins fail.
The total \(g\) of qualifying pins satisfies
\(g\ge s+u-2t-42=3m-2n-2s-42\ge3m-2n-2g-42\).
This proves the one-circle estimate. In the coverage setting it handles
coinciding circles or a layer with fewer than \(0.07n\) points, since the
other then has more than \(0.92n\) points.

If the two concentric circles differ in radius and both layers are large,
their rotations have the same group \(H=\mu_h\). Normalize the direction
of \(p_0\) to one. Each of the nearly full orbits contains a good pin
because \(0.89h>0.111n\). Its phase \(\alpha\) therefore satisfies
\(\alpha^2\in H\), so the two core orbits lie in \(\mu_{2h}\). Each core
surrounds the center, and strict convexity allows at most one point per
ray, so the orbits are the two distinct alternating cosets.
A pin outside this grid has all within-core distances distinct on each
core: its reflection does not preserve either orbit. Mixed coincidences
between the two cores at that pin are at most 44, by the same five-term
Laurent equation and Newton polygon area two as in the majority argument.
It consequently sees more than \(1.78h-44>n/2\) distances for sufficiently
large \(n\). Discard the \(k\le0.01n\) off-circle targets and \(\ell<n/10\)
off-grid targets in the two layers for an initial reduced set, of size
\(n-K\) with \(K=k+\ell\le0.11n\).
The two cores have fewer than \(0.11h\) gaps each. There is an inner core
vertex with its two outer neighbors present, since the omissions exclude
fewer than \(0.33h\) positions. This vertex must lie outside the chord
of the neighbors, giving radius ratio \(\rho>\cos(\pi/h)\).

Put \(\tau=\pi/h\) and scale the larger radius to one. From a pin of
radius \(a\in[\rho,1]\), the angular bin \(j\tau\) gives distances
strictly increasing in target radius. Bins \(j,j+1\) are separated for
\(1\le j\le h-2\): at \(\rho=\cos\tau\), the difference between the
minimum squared distance in the later bin and maximum in the earlier
one equals \(2a\sin\tau\sin((j+1)\tau)-\sin^2\tau>0\).
This difference increases with \(\rho\) on the relevant interval.
An outer point without antipode therefore qualifies within the reduced set.
If it has an antipode, that antipode qualifies because the last bin has the
outer point alone, with the strictly greatest distance. Taking each outer
point or its antipode by this rule gives distinct pins, more than \(0.89h\),
each with at least \(\lfloor(n-K)/2\rfloor\) reduced-set distances.

On adding the discarded targets, a new target can merge with the old
two layers at at most 176 selected pins. For unequal target radii, for
each of two pin radii and two target radii the five-term equation in
the pin direction and target relative direction has at most 44 torsion
solutions. An off-grid target in one layer cannot merge with a grid target
in its own layer because reflection at a grid pin preserves the grid.
Pairs of additional targets create at most two equalities each over all
pins. Hence total losses are at most \(K(K+175)\). When \(K>0\), at most
\(2(K+175)\) pins can lose more than \(K/2\). All other selected pins reach
\(\lfloor(n-K)/2\rfloor+K-\lfloor K/2\rfloor\ge\lfloor n/2\rfloor\).
When \(K=0\) they already reach that bound. For large \(n\), the count
remaining exceeds \(0.89(0.4445n)-0.22n-350>n/10\), a contradiction.
\end{proof}

\begin{proposition}[982-M03, 982-M04]
The union of two complete concentric regular polygons of orders at least
three reaches the half bound whenever it is in strictly convex position.
For arbitrary samples on two concentric circles, every pin reaches the
half bound if the radius ratio satisfies no nonzero polynomial of degree
at most two over the field generated by all cosines of sample direction differences.
\end{proposition}
\begin{proof}
The equal-radius case is cocircular. Otherwise denote the complete
polygons by \(A,B\), of orders \(m\ge k\), and scale the radius of \(A\) to
one. If some pin of \(B\) lies off all symmetry axes of \(A\), it sees the
\(m\) distances to \(A\) distinctly. Otherwise no same-ray pair is possible
by strict convexity, and all directions of \(B\) bisect adjacent directions
of \(A\). Thus \(H=m/k\) is an integer. Put \(h=\pi/m\) and \(f(s)=2\sin(s/2)\).
All cross angles are odd multiples of \(h\); \(A\)'s own chord angles
are even multiples. Write \(R\) for the radius of \(B\).

If \(R<1\), strict convexity at any \(B\) vertex between two \(A\) directions
gives \(R>\cos h\). A cross distance at angle \(s\) is below \(f(s)\)
by the cosine rule and less than \(1-\cos h\) below it by the triangle
inequality. For \(s\le\pi-h\),
\[
 f(s)-f(s-h)\ge f(\pi-h)-f(\pi-2h)>1-\cos h.
\]
The final difference after subtracting \(1-\cos h\) equals
\(4\sin^2(h/4)\cos(h/2)\). Thus only the antipodal cross angle, possible
for odd \(m\), may collide with an \(A\)-chord distance.

If \(H=1\), exchange the roles to make \(A\) outer. At an outer pin
the within-layer and cross bins contribute \(\lfloor m/2\rfloor\) and
\(\lceil m/2\rceil\) distances before merging. If the antipodal distance
merges, it equals \(f(\pi-h)=2\cos(h/2)\), because \(1+R>1+\cos h>
2\cos(3h/2)\), the second largest candidate (or a vacuous candidate
when \(m=3\)). The second inequality follows from
\(c^2-\cos(3h/2)=c(1-c)(4c+3)>0\) for \(c=\cos(h/2)\).
So \(R=2c-1\). At an inner pin rescale the inner radius to one.
The cross distances for \(s\le\pi-2h\) are then strictly above \(f(s)\)
and at most \(1/R-1\) higher. The next possible chord distance is separated
by at least \(f(\pi-h)-f(\pi-2h)=2(1-c)(2c+1)>1/R-1\), since
\((2c+1)(2c-1)>1\). The antipodal cross distance exceeds two.
The inner pin therefore sees \(m\) distances.

For \(H\ge2\), suppose first \(m\ge9\). We claim at most four cross
targets at any \(B\) pin have lengths in the old \(A\)-chord set.
The inner-\(B\) case already gives at most one. If \(R>1\), rotate
the \(B\) pin to the positive horizontal ray. The orientation of it
and the \(A\) vertices at angles \(h,3h\) forces \(R<\cos h/\cos(2h)\):
the chord of those vertices has inward half-plane
\(x\cos(2h)+y\sin(2h)<\cos h\). A cross distance is above \(f(s)\)
and at most \(R-1\) higher. For \(s\le\pi-5h\) the next possible chord
is farther than this, since
\[
 f(\pi-4h)-f(\pi-5h)=4\sin(9h/4)\sin(h/4)>
 \cos h/\cos(2h)-1 .
\]
After cancellation the inequality reads
\(\sin(9h/4)\cos(2h)>\sin(3h/2)\cos(h/4)\).
For \(0<h\le\pi/9\), the sine chord bound on \([0,\pi/4]\) gives
\(\sin(9h/4)>2h\), and \(\cos(2h)>3/4\). The right hand side is at most
\(3h/2\). Only two nonantipodal odd bins remain above \(\pi-5h\) in
each parity, and the antipodal distance when present exceeds two.
There are at most two targets in each bin, proving the claim.

At a pin of \(A\) off the axes of \(B\), the \(k\) cross distances
are distinct, so failure would require at least \(\lceil k/2\rceil\)
collisions with its \(A\) chord distances. The axis congruence is
\(H\mid 2j-1\) for \(A\) direction \(2jh\). There are no on-axis
pins for even \(H\); all pins failing would require
\(m\lceil k/2\rceil>4k\), a contradiction. For odd \(H\ge3\) there
are \(k\) on-axis pins, each with \(\lceil k/2\rceil\) distinct
cross distances by reflection. The two orders have the same parity,
so each on-axis pin must suffer at least one merge to fail.
The combined requirement \((m-k)\lceil k/2\rceil+k\ge5k\)
contradicts the four-per-\(B\)-pin bound.

For \(m<9\) the proper-divisor cases are \((m,k)=(6,3),(8,4)\).
For \((6,3)\) failure would make the three cross distances at an \(A\)
pin precisely the \(A\) chords, with squares \(1,3,4\).
Their sum \(3(1+R^2)=8\) would give a right-angle cross square
\(1+R^2=8/3\), impossible. For \((8,4)\) the two layers together
are centrally symmetric and lie in the minimum disk with an
outer contact diameter, so the support-chain lemma applies.

For the arbitrary-sample claim, scale the radii to \(1,R\).
A mixed coincidence at a pin on the first or second layer,
respectively, implies
\[
 2-2\cos u=1+R^2-2R\cos v,\qquad
 R^2(2-2\cos u)=1+R^2-2R\cos v,
\]
with \(u\) the same-layer angle. Each is a nonzero degree-at-most-two
polynomial over the specified field, excluded by hypothesis. Each
separate layer has distance multiplicity at most two, and without
mixed coincidences the two contributions give the half bound.
\end{proof}

\section{An incomplete odd regular core}

\begin{theorem}[982-R07, 982-R08]
Let \(A\) consist of \(m\) vertices of a unit regular odd \(M\)-gon, and \(B\)
consist of \(t\) points off the unit circle, with \(A\cup B\) strictly convex.
If \(M\ge11\) and \(\cos h<|b|\le\cos h/\cos(2h)\) for all \(b\in B\),
where \(h=\pi/M\), \(m\ge3t/2+83/2\) guarantees a qualifying pin in \(A\).
If \(B\) is cocircular, then \(M\ge3\) and \(m\ge t+44\) suffice without the annular hypothesis.
\end{theorem}

\begin{proof}
First, any three distinct lines spanned by vertices of an odd regular polygon
that meet at a finite point meet at a polygon vertex. For concurrence at
a nonvertex the endpoints are all distinct. Normalize the three pairs to
\((1,r),(a,b),(c,d)\) with entries in \(\mu_M\). The equation of the second
line is \(z+ab\bar z=a+b\). Substituting \(z=1+\lambda(r-1)\) with real
\(\lambda\), and comparing the third line, gives
\[
 (1-a)(1-b)(1-cd/r)=(1-c)(1-d)(1-ab/r)\ne0.
\]
The nonzero assertion uses nonparallelism and distinct endpoints.
Apply the cyclotomic field automorphism squaring the roots, then divide
by this equality. The resulting equality of products with plus signs,
together with the common product \(abcd/r\), forces equality of the
first two elementary symmetric coefficients. Thus \((a,b,cd/r)\)
and \((c,d,ab/r)\) are the same multiset. Since \(a\notin\{c,d\}\) this
forces \(b=r\), a contradiction.

For the annular case put \(S=\sin h,\ C=\cos h,\ U=C/\cos(2h)\).
Form a graph on \(B\) by joining pairs whose bisector meets two vertices
of \(A\). A triangle of this graph would have its circumcenter at
a vertex of \(A\), by the preceding concurrence property. Rotate this
common vertex to 1 and write the three targets as \(1+R e^{i\theta_i}\).
They lie in one open half-plane at 1, as do all other vertices, so
choose \(\theta_1<\theta_2<\theta_3<\theta_1+\pi\). The second unit-circle
intersection of the \(i,j\) bisector is
\(-e^{i(\theta_i+\theta_j)}\in A\), lying on the positive ray at angle
\((\theta_i+\theta_j)/2\) from 1. For pairs \(12,23\) these rays enclose
\(\theta_2\), with angular gap less than \(\pi/2\). Both unit-circle
intersections are at least \(2S\) from 1. Their joining segment has
projection strictly greater than \(\sqrt 2 S\) along the mid-direction.
Since the middle target must be outside the triangle bounded by this
segment and the vertex 1, \(R>\sqrt2 S\).

From the three second intersections being roots of unity, the
\(\theta_i\) are congruent modulo \(2h\) and their pairwise sums are
\(\pi\) modulo \(2h\). The folded angles
\(\gamma_i=\arccos(\cos\theta_i)\) are distinct half-integer multiples
of \(h\) in \([h/2,\pi-h/2]\). In fact two original angles giving
the same folded angle would either coincide or sum to zero modulo
\(2\pi\), incompatible with the pairwise sum congruence for odd \(M\).
Thus the maximum and minimum folded angles differ by at least \(2h\).

The allowed squared-radius width is
\[
 U^2-C^2=4S^2\left(\frac{1-S^2}{1-2S^2}\right)^2<\frac{121}{25}S^2
\]
since \(h\le\pi/11<2/7\) gives \(S^2<1/12\). The actual width equals
\(2R\) times the span of the cosines of \(\gamma_i\). For \(R\ge1\)
the least span over any interval of length \(2h\) in
\([h/2,\pi-h/2]\) is \(2\sin(3h/2)S\), and
\(\sin(3h/2)/S>3\sqrt3/4>121/100\). This already exceeds the allowed width.
For \(\sqrt2 S<R\le1\), the radial bounds give
\[
 -\frac{R^2+S^2}{2R}<\cos\gamma_i\le\frac{U^2-1}{2R}<\frac{54S^2}{25R}.
\]
Here \(U^2-1=(3S^2-4S^4)/(1-2S^2)^2\).
If \(R\ge3S/2\), both extreme cosine magnitudes are bounded
by \(13/24\) using \(S^2<1/12\) and monotonicity of \(R+S^2/R\).
Then the sine between the extreme angles exceeds \(5/6\), so their
squared-radius span exceeds \(5Sh>121S^2/25\). If \(R\le3S/2\),
both magnitudes are below \(1/2\) using \(\sqrt2 S<R\), so the width
exceeds \(2\sqrt6 Sh>121S^2/25\). This contradicts the annulus in
all cases. The graph is triangle-free, hence has at most
\(\lfloor t^2/4\rfloor\) edges by Mantel's inequality.

Each nonzero \(b\in B\) causes at most 22 merges with old circle
distances over pins in \(A\). The five-term argument for the majority
theorem bounds the number of oriented torsion solutions by 44, and
each merge supplies distinct \(u,u^{-1}\): \(u=1\) would place \(b\)
on the circle, and \(u=-1\) is excluded by odd \(M\).
If the origin belongs to \(B\), strict support there confines \(A\)
to an open semicircle and an extreme pin of \(A\) suffices by the
threshold. Hence we may exclude that case. In the annular case the
total double-witness count from pairs of new targets is bounded by
\(\binom t2+\lfloor t^2/4\rfloor\), counting at most one per pair and
one extra per graph edge. Failure at all core pins would require
\[
 m(t+1)/2\le22t+\binom t2+\lfloor t^2/4\rfloor
 \le 3t^2/4+43t/2.
\]
This contradicts \(m\ge3t/2+83/2\).

For the cocircular case let \(Q\) be the center of \(B\). If \(Q\notin A\),
all bisectors of \(B\)-pairs pass through \(Q\), with at most two lines
meeting \(A\) twice by the concurrence result. If \(Q\) is itself an
unselected polygon vertex, no such line can meet \(A\) twice. Each line
accounts for at most \(\lfloor t/2\rfloor\) disjoint reflected pairs.
Consequently total pair merges are at most \(\binom t2+2\lfloor t/2\rfloor\).
With the 22-target penalty the failure condition gives
\(m(t+1)\le t^2+45t\), impossible for \(m\ge t+44\).
If \(Q\in A\), omit its pin but keep all original \(A\)-targets.
Each pair of new targets then accounts for at most one merge, since
one of its two circle intersections is \(Q\). Thus failure requires
\((m-1)(t+1)\le44t+2\binom t2=t^2+43t\), also impossible.
\end{proof}


\begin{thebibliography}{9}
\bibitem{erdos} P. Erd\H{o}s, On sets of distances of \(n\) points, \emph{Amer. Math. Monthly} 53 (1946), 248--250.
\href{https://www.renyi.hu/~p_erdos/1946-03.pdf}{Author archive}.
\bibitem{moser} L. Moser, On the different distances determined by \(n\) points, \emph{Amer. Math. Monthly} 59 (1952), 85--91.
\href{https://doi.org/10.1080/00029890.1952.11988075}{DOI}.
\bibitem{ef} P. Erd\H{o}s and P. Fishburn, A postscript on distances in convex \(n\)-gons, \emph{Discrete Comput. Geom.} 11 (1994), 111--117.
\href{https://doi.org/10.1007/BF02573998}{DOI}.
\bibitem{nppz} G. Nivasch, J. Pach, R. Pinchasi, and S. Zerbib, The number of distinct distances from a vertex of a convex polygon, \emph{J. Comput. Geom.} 4 (2013), 1--12.
\href{https://arxiv.org/abs/1207.1266}{arXiv:1207.1266}.
\bibitem{kominers} S. D. Kominers, A (slightly) stronger lower bound on the number of distinct distances from a vertex of a convex polygon, manuscript, 25 July 2026.
\href{https://scottkom.com/assets/articles/Kominers_Distinct_Distances_from_a_Vertex.pdf}{Author manuscript}.
\bibitem{st} E. Szemer\'edi and W. T. Trotter, Extremal problems in discrete geometry, \emph{Combinatorica} 3 (1983), 381--392.
\href{https://trotter.math.gatech.edu/papers/38.pdf}{Author copy}.
\bibitem{sz} J. Solymosi and J. Zahl, Improved Elekes--Szab\'o type estimates using proximity, 2022.
\href{https://arxiv.org/abs/2211.13294}{arXiv:2211.13294}.
\bibitem{ess} J.-H. Evertse, H. P. Schlickewei, and W. M. Schmidt, Linear equations in variables which lie in a multiplicative group, \emph{Ann. of Math.} 155 (2002), 807--836. \href{https://annals.math.princeton.edu/2002/155-3/p04}{Article}.
\bibitem{bs} F. Beukers and C. J. Smyth, Cyclotomic points on curves. \href{https://webhomes.maths.ed.ac.uk/~chris/preprints/beukers_smyth.pdf}{Author copy}.
\bibitem{msw} H. Martini, K. J. Swanepoel, and G. Wei, The geometry of Minkowski spaces: a survey, Part I. \href{https://arxiv.org/abs/0708.2900}{Author manuscript}.
\bibitem{dgos} Z. Dvir, A. Garg, R. Oliveira, and J. Solymosi, Rank bounds for design matrices with block entries and geometric applications. \href{https://arxiv.org/abs/1610.08923}{arXiv:1610.08923}.
\bibitem{heintz} J. Heintz, Definability and fast quantifier elimination in algebraically closed fields, \emph{Theoret. Comput. Sci.} 24 (1983), 239--277. \href{https://doi.org/10.1016/0304-3975(83)90002-6}{DOI}.
\end{thebibliography}
\end{document}
